\begin{topic}[id=lorawan_ttn,
             difficulty={Einstieg},
             type={Protokoll + Praxisbezug},
             prerequisites={Grundlagen Rechnernetze; Basiswissen Funkkommunikation hilfreich},
             practice={gut möglich},
             owner={Dozententeam},
             updated={2026-04-09},
             tags={ LoRaWAN, TTN, Gateways, ADR, Klassen A/B/C, Payload, Security }]{LoRaWAN \& The Things Network}

\TopicDescription{LoRaWAN ermöglicht energiearme Weitverkehrskommunikation mit sehr kleinen Datenraten – ideal für Sensorik über Kilometer. Du lernst Geräteklassen (A/B/C), Adaptive Data Rate (ADR), Join/Session-Keys und typische Payload-Designs. The Things Network (TTN) dient als offenes Ökosystem für schnelle Tests. Wir besprechen Reichweite vs. Airtime-Politik, Duty-Cycle-Regeln, Sicherheitsaspekte (Keys, Frame Counter) und realistische Use-Cases. Ziel ist eine fundierte Einschätzung, wo LoRaWAN glänzt und wo andere LPWANs (NB-IoT) überlegen sind.}

\TopicLeitfragen{\item Welche Grenzen setzt der Duty-Cycle praktisch?
\item Wie gestaltet man eine robuste Payload/Decoding?
\item Wann ist NB-IoT statt LoRaWAN sinnvoller?
}
\TopicScope{\item Kein Hardware-Gateway-Bau.
\item Keine tiefen Regulierungsdetails.
}
\TopicPractice{Optionale Praxisidee: TTN-Endgerät registrieren, Payload senden, einfache Visualisierung.}

\TopicKeywords{LoRaWAN, TTN, Gateways, ADR, Klassen A/B/C, Payload, Security}

\nocite{lorawanSpec,ttnDocs}
\end{topic}
